\documentclass[10pt]{article}

% Packages pour les marges
\usepackage[
    top=2cm,
    bottom=2cm,
    left=2cm,
    right=2cm
]{geometry}

% Police sans serif
% \usepackage{helvet}
% \renewcommand{\familydefault}{\sfdefault}

% Packages existants
\usepackage[french]{babel}
\usepackage[utf8]{inputenc}
\usepackage[T1]{fontenc}
\usepackage{amsmath}
\usepackage{amsfonts}
\usepackage{amssymb}
\usepackage[version=4]{mhchem}
\usepackage{stmaryrd}
\usepackage{listings}

\usepackage{enumitem}
\usepackage{ifthen}
\usepackage{graphicx}
\usepackage{verbatim}
\usepackage{url}
\usepackage{mdframed}
\usepackage[sf]{titlesec}
\titleformat{\section}{\sffamily\large\bfseries}{\thesection}{1em}{}
% Définition de la variable pour afficher les corrections
\newboolean{showSolutions}
% Décommentez la ligne suivante pour afficher les solutions
\input \jobname.adr
% \input \jobname.adr


\title{TD3 - Analyse Numérique}

\author{}
\date{}

\newenvironment{solution}
    {\par\vspace{0.5em}\begin{mdframed}[linewidth=0.5pt]\par}
    {\end{mdframed}\par\vspace{0.5em}}

\begin{document}
\sffamily
\maketitle

\section*{Des polynômes familiers}
On note $(p_n)_{n\in\mathbb{N}}$ la suite des nombres premiers. 
\begin{enumerate}
    \item Donner le polynôme de degré minimal tel que pour chaque $k\in\{1, \cdots, n\}$ $P(k)$ soit le $k$-ième nombre premier. 
    \ifthenelse{\boolean{showSolutions}}{
        \begin{solution}
            On commence par former les polynômes de Lagrange $\ell_0, \cdots, \ell_n$ tels que $\ell_i(j) = \delta_{ij}$:
            \[
            \ell_i(x) = \prod_{\substack{0 \leq j \leq n \\ j \neq i}} \frac{x-j}{i-j}
            \]
            On peut alors écrire $P(x) = \sum_{i=0}^n p_i \ell_i(x)$.
        \end{solution}
    }{}
    \item A votre avis, $P(n+1)$ est-il premier ?
    \ifthenelse{\boolean{showSolutions}}{
        \begin{solution}
            Aucune chance...  
            Il n'y a même aucune raison pour que $P(n+1)$ soit un nombre entier. 
        \end{solution}
    }{}
\end{enumerate}


\vspace{1cm}
\section*{Deux décompositions LU}


\begin{enumerate}
    \item On donne la matrice $
    A=\left(\begin{array}{cccc}
    2 & -1 & 0 & 0 \\
    -1 & 2 & -1 & 0 \\
    0 & -1 & 2 & -1 \\
    0 & 0 & -1 & 2
    \end{array}\right)
    $
    Former la décomposition $L U$ de $A$, où $L$ est une matrice triangulaire inférieure dont les éléments diagonaux sont égaux à l'unité et $U$ est une matrice triangulaire supérieure. 
    \ifthenelse{\boolean{showSolutions}}{
        \begin{solution}
            $A = LU$ avec 
            \[
            L = \left(\begin{array}{cccc}
                    1 & 0 & 0 & 0 \\ 
                    -1/2 & 1 & 0 & 0 \\ 
                    0 & -2/3 & 1 & 0 \\ 
                    0 & 0 & -3/4 & 1
                \end{array}\right), \qquad 
            U = \left(\begin{array}{cccc}
                    2 & -1 & 0 & 0 \\
                    0 & 3/2 & -1 & 0 \\
                    0 & 0 & 4/3 & -1 \\
                    0 & 0 & 0 & 5/4
                \end{array}\right)
            \]
        \end{solution}
    }{}
    
    \item Calculer $D=\operatorname{det} A$.
    \ifthenelse{\boolean{showSolutions}}{
        \begin{solution}
            $D = \operatorname{det} A = \operatorname{det} L \cdot  \operatorname{det} U = 1 \cdot \frac{3}{2} \cdot \frac{4}{3} \cdot \frac{5}{4} = 5$
        \end{solution}
    }{}
    \item Utiliser la factorisation précédente pour résoudre le système linéaire $A x = b$, avec $b = [2,3,4,2]$.
    \ifthenelse{\boolean{showSolutions}}{
        \begin{solution}
            On cherche à résoudre $LUX = b$, on pose $Y = UX$ et on résout $LY = b$. Comme on connait $L^{-1}$, on obtient directement 
            \[
            Y = L^{-1} b = \left(\begin{array}{cccc}
                1 & 0 & 0 & 0 \\ 
                1/2 & 1 & 0 & 0 \\ 
                1/3 & 2/3 & 1 & 0 \\ 
                1/4 & 1/2 & 3/4 & 1
            \end{array}\right) \left(\begin{array}{c}
                2 \\
                3 \\
                4 \\
                2
            \end{array}\right) = \left(\begin{array}{c}
                2 \\
                4 \\
                20/3 \\
                7
            \end{array}\right)
            \]


            On résout ensuite $UX = Y$, le système linéaire s'écrit : 
            \[
            \left(\begin{array}{cccc}
                2 & -1 & 0 & 0 \\
                0 & 3/2 & -1 & 0 \\
                0 & 0 & 4/3 & -1 \\
                0 & 0 & 0 & 5/4
            \end{array}\right) \left(\begin{array}{c}   
                x_1 \\
                x_2 \\
                x_3 \\
                x_4
            \end{array}\right) = \left(\begin{array}{c}
                2 \\
                4 \\
                20/3 \\
                7
            \end{array}\right)
            \]

            On trouve $x_4 = 28/5$, $x_3 = 46/5$, $x_2 = 44/5$ et $x_1 = 54/10 = 27/5$.
            
            En vérifiant, on obtient bien $A x = b$.
        \end{solution}
    }{}


\item Ecrire la décomposition LU de la matrice 
$\left(\begin{array}{ccc}1 & 2 & -3 \\ 2 & 6 & -5 \\ 1 & -2 & 7\end{array}\right)$
\ifthenelse{\boolean{showSolutions}}{
        \begin{solution}
            $A = LU$ avec 
            $L = \left(\begin{array}{ccc}1 & 0 & 0 \\ 2 & 1 & 0 \\ 1 & -2 & 1\end{array}\right)$
            $U = \left(\begin{array}{ccc}1 & 2 & -3 \\ 0 & 2 & 1 \\ 0 & 0 & 12\end{array}\right)$
                    \end{solution}
    }{}

\item Résoudre le système linéaire 
$\left\{\begin{aligned} x_1+2 x_2-3 x_3 & =1 \\ 2 x_1+6 x_2-5 x_3 & =-1 \\ x_1-2 x_2+7 x_3 & =2\end{aligned}\right.$

\ifthenelse{\boolean{showSolutions}}{
        \begin{solution}
            On résout par descente $Ly=b$ et on trouve $y=(1,-3,-5)^t$. 
            
            Puis on résout $Ux=y$ par remontée, et on trouve $x=\left(\frac{7}{3},-\frac{31}{24},-\frac{5}{12}\right)^t$.
        \end{solution}
    }{}
    
    \item En raisonnant par l'absurde, montrer que la décomposition LU est unique, si $\boldsymbol{L}$ une matrice triangulaire inférieure avec des $1$ sur la diagonale et $\boldsymbol{U}$ une matrice triangulaire supérieure. 
    \ifthenelse{\boolean{showSolutions}}{
        \begin{solution}
            Supposons que $A = LU = L'U'$ avec $L$ et $L'$ triangulaires inférieures avec des $1$ sur la diagonale et $U$ et $U'$ triangulaires supérieures. 
            $L$ et $L'$ sont donc inversibles. Si $U$ ou $U'$ n'est pas inversible, alors $A$ ne serait pas inversible non plus. 
            
            Alors $L'^{-1}L = UU'^{-1}$. Or $L'^{-1}L$ est triangulaire inférieure avec des $1$ sur la diagonale et $UU'^{-1}$ est triangulaire supérieure, donc $L'^{-1}L = I$. Donc $L = L'$ et $U = U'$.
        \end{solution}
    }{}
    \vspace{1em}

    On peut obtenir les matrices $L$ et $U$ à partir de l'égalité $A = LU$ sans effectuer l'algorithme de Gauss:
    
    \item En posant des inconnues pour les matrices $L$ et $U$ de taille $3\times 3$, déterminer les équations vérifiées par les coefficients de $L$ et $U$.\newline
    On pourra regarder dans l'ordre la première ligne de $A = LU$, puis la première colonne, ensuite la deuxième ligne puis la deuxième colonne, ... 
    \ifthenelse{\boolean{showSolutions}}{
        \begin{solution}
            Pour des matrices $3\times 3$: 

            On pose $L = \left(\begin{array}{ccc}
                1 & 0 & 0 \\
                l_{21} & 1 & 0 \\
                l_{31} & l_{32} & 1
            \end{array}\right)$ et $U = \left(\begin{array}{ccc}
                u_{11} & u_{12} & u_{13} \\
                0 & u_{22} & u_{23} \\ 
                0 & 0 & u_{33}
            \end{array}\right)$.

            Regardons la première ligne de l'équation $A = LU$ : 
            \[
            \left(\begin{array}{ccc}
                a_{11} & a_{12} & a_{13} \\
            \end{array}\right) = \left(\begin{array}{ccc}
                u_{11} & u_{12} & u_{13} \\
            \end{array}\right)
            \]
            On obtient donc directement la première ligne de $U$. 

            Regardons ensuite la première coolonne de l'équation $A = LU$ : 
            \[
            \left(\begin{array}{c}
                a_{11} \\
                a_{21} \\
                a_{31}
            \end{array}\right) = \left(\begin{array}{c}
                u_{11} \\
                l_{21}u_{11} \\
                l_{31}u_{11}
            \end{array}\right)
            \]
            On obtient donc facilement la première colonne de $L$: 
            \[
            l_{21} = \frac{a_{21}}{a_{11}}, \quad l_{31} = \frac{a_{31}}{a_{11}}
            \]

            Regardons ensuite la deuxième ligne de l'équation $A = LU$ : 
            \[
            \left(\begin{array}{ccc}
                a_{21} & a_{22} & a_{23} \\
            \end{array}\right) = \left(\begin{array}{ccc}
                l_{21}u_{11} & l_{21}u_{12} + u_{22} & l_{21}u_{13} + u_{23} \\
            \end{array}\right)
            \]
            On peut donc en déduire la deuxième ligne de $U$ ...            
            
        \end{solution}
    }{}
    
    \item Ecrire un algorithme qui calcule $L$ et $U$ à partir de $A$ sans effectuer l'algorithme de Gauss.


    \end{enumerate}
    \vspace{1em}

    \section*{Méthode de Jacobi}
    On donne les deux systèmes linéaires :
    
    $$
    \left[\begin{array}{cc}
    10 & 1 \\
    2 & 10
    \end{array}\right]\left[\begin{array}{l}
    x_1 \\
    x_2
    \end{array}\right]=\left[\begin{array}{l}
    11 \\
    12
    \end{array}\right]; \qquad\left[\begin{array}{cc}
    1 & 10 \\
    10 & 2
    \end{array}\right]\left[\begin{array}{l}
    x_1 \\
    x_2
    \end{array}\right]=\left[\begin{array}{l}
    11 \\
    12
    \end{array}\right]
    $$
    \begin{enumerate}
        \item Trouver les solutions exactes.
        \ifthenelse{\boolean{showSolutions}}{
            \begin{solution}
                On peut utiliser un pivot de Gauss, pour le premier système, changer $L_2$ en $L_2 - L_1/5$ donne
                \[
                \left(\begin{array}{cc}
                    10 & 1 \\
                    0 & 49/5 \\
                \end{array}\right) \left(\begin{array}{l}
                    x_1 \\
                    x_2
                \end{array}\right) = \left(\begin{array}{l}11 \\ 49/5\end{array}\right)
                \]
                On peut alors résoudre le système pour trouver $x_1 = 1$ et $x_2 = 1$.

                Le second système admet la même solution. 
            \end{solution}
        }{}
        \item Résoudre chaque système par la méthode itérative de Jacobi, à partir du vecteur initial $x^{(0)}=(0,0)^T$, jusqu'au vecteur $x^{(4)}$ inclus.
        \ifthenelse{\boolean{showSolutions}}{
            \begin{solution}
            La méthode de Jacobi est une méthode itérative pour résoudre les systèmes linéaires de la forme $Ax = b$. 
            Elle repose sur la décomposition de la matrice $A$ en trois matrices : $A = D - E - F$, où $D$ est la matrice diagonale contenant les éléments diagonaux de $A$, $E$ est la matrice triangulaire inférieure contenant les éléments sous la diagonale, et $F$ est la matrice triangulaire supérieure contenant les éléments au-dessus de la diagonale.

            \begin{itemize}
                \item[$\bullet$] D'abord on choisit un vecteur initial, dans notre cas, nous avons $x^{(0)} = (0,0)^T$.

                \item[$\bullet$] On calcule le vecteur $x^{(k+1)}$ à partir de $x^{(k)}$ en utilisant la formule matricielle suivante :
                \[
                x^{(k+1)} = D^{-1}(b + Ex^{(k)} + Fx^{(k)})
                \]
                où $D^{-1}$ est l'inverse de la matrice diagonale $D$.

                \item[$\bullet$] On répète l'étape 2 jusqu'à ce que la solution converge, c'est-à-dire jusqu'à ce que la différence entre deux itérations successives soit inférieure à un certain seuil de tolérance, ou jusqu'à atteindre un nombre maximum d'itérations.
            \end{itemize}

            Appliquons cette méthode aux deux systèmes donnés :

            \[
            \left[\begin{array}{cc}
            10 & 1 \\
            2 & 10
            \end{array}\right]\left[\begin{array}{l}
            x_1 \\
            x_2
            \end{array}\right]=\left[\begin{array}{l}
            11 \\
            12
            \end{array}\right]
            \]
            \begin{itemize}
                \item[$\bullet$] Décomposition : $D = \left[\begin{array}{cc} 10 & 0 \\ 0 & 10 \end{array}\right]$, $E = \left[\begin{array}{cc} 0 & 0 \\ -2 & 0 \end{array}\right]$, $F = \left[\begin{array}{cc} 0 & -1 \\ 0 & 0 \end{array}\right]$
                \item[$\bullet$] Itération 0 : $x^{(0)} = (0,0)^T$
                \item[$\bullet$] Itération 1 : $x^{(1)} = D^{-1}(b + Ex^{(0)} + Fx^{(0)}) = \left[\begin{array}{c} 1.1 \\ 1.2 \end{array}\right]$
                \item[$\bullet$] Itération 2 : $x^{(2)} = D^{-1}(b + Ex^{(1)} + Fx^{(1)}) = \left[\begin{array}{c} 0.98 \\ 0.98 \end{array}\right]$
                \item[$\bullet$] Itération 3 : $x^{(3)} = D^{-1}(b + Ex^{(2)} + Fx^{(2)}) = \left[\begin{array}{c} 1.002 \\ 1.004 \end{array}\right]$
                \item[$\bullet$] Itération 4 : $x^{(4)} = D^{-1}(b + Ex^{(3)} + Fx^{(3)}) = \left[\begin{array}{c} 0.9996 \\ 0.9996 \end{array}\right]$
            \end{itemize}

            **Deuxième système :**
            \[
            \left[\begin{array}{cc}
            1 & 10 \\
            10 & 2
            \end{array}\right]\left[\begin{array}{l}
            x_1 \\
            x_2
            \end{array}\right]=\left[\begin{array}{l}
            11 \\
            12
            \end{array}\right]
            \]

            \begin{itemize}
                \item[$\bullet$] Décomposition : $D = \left[\begin{array}{cc} 1 & 0 \\ 0 & 2 \end{array}\right]$, $E = \left[\begin{array}{cc} 0 & 0 \\ -10 & 0 \end{array}\right]$, $F = \left[\begin{array}{cc} 0 & -10 \\ 0 & 0 \end{array}\right]$
                \item[$\bullet$] Itération 0 : $x^{(0)} = (0,0)^T$
                \item[$\bullet$] Itération 1 : $x^{(1)} = D^{-1}(b + Ex^{(0)} + Fx^{(0)}) = \left[\begin{array}{c} 11 \\ 6 \end{array}\right]$
                \item[$\bullet$] Itération 2 : $x^{(2)} = D^{-1}(b + Ex^{(1)} + Fx^{(1)}) = \left[\begin{array}{c} -49 \\ -49 \end{array}\right]$
                \item[$\bullet$] Itération 3 : $x^{(3)} = D^{-1}(b + Ex^{(2)} + Fx^{(2)}) = \left[\begin{array}{c} 501 \\ 251 \end{array}\right]$
                \item[$\bullet$] Itération 4 : $x^{(4)} = D^{-1}(b + Ex^{(3)} + Fx^{(3)}) = \left[\begin{array}{c} -2499 \\ -2499 \end{array}\right]$
            \end{itemize}

            \end{solution}
        }{}
        \item Expliquer qualitativement le comportement différent des deux suites de vecteurs.
        \ifthenelse{\boolean{showSolutions}}{
            \begin{solution}
                On observe que pour le premier système, la méthode de Jacobi converge rapidement vers la solution exacte $(1, 1)^T$. En revanche, pour le deuxième système, la méthode semble diverger.
            \end{solution}
        }{}
    \end{enumerate}
    
    \ifthenelse{\boolean{showSolutions}}{}
    {
       \newpage 
    }
    

    \vspace{1cm}
    \section*{Une décomposition LU qui se passe mal}

    \begin{enumerate}
        \item Montrer que la matrice 
$
A=\left(\begin{array}{lll}
0 & 3 & 0 \\
1 & 2 & 0 \\
3 & 5 & 2
\end{array}\right)$
        n'est pas décomposable sous la forme $L U$
        \ifthenelse{\boolean{showSolutions}}{
            \begin{solution}
                Le premier pivot est nul donc nous ne pouvons pas effectuer la décomposition $LU$.
            \end{solution}
        }{}
        \item Trouver $P$ une matrice de permutation et $B$ une matrice décomposable sous la forme $L U$ telles que $B=P A$ et déduire une factorisation $L U$ de la matrice $B$.
        \ifthenelse{\boolean{showSolutions}}{
            \begin{solution}
                On peut permuter les lignes 1 et 2 pour obtenir la matrice $B$ :
                \[
                B = \left(\begin{array}{ccc}
                1 & 2 & 0 \\    
                0 & 3 & 0 \\
                3 & 5 & 2
                \end{array}\right)
                \]
                Cela correspond à multiplier $A$ à gauche par la matrice de permutation $P$ :
                \[
                P = \left(\begin{array}{ccc}
                0 & 1 & 0 \\
                1 & 0 & 0 \\
                0 & 0 & 1
                \end{array}\right)
                \]
                On peut alors effectuer la décomposition $LU$ de $B$. 
                \end{solution}
        }{}
        
        \item Utiliser ce résultat pour résoudre $A\left(\begin{array}{l}x \\ y \\ z\end{array}\right)=\left(\begin{array}{l}1 \\ 0 \\ 1\end{array}\right)$.
\end{enumerate}

\vspace{1cm}

\section*{Vers la résolution des équations différentielles}
Seule une infime minorité des équations différentielles est soluble mathématiquement.

% Discrétisation, intégration sur chaque intervalle. Approximation de l'intégrale. 
% Runge Kutta, utilisation des formules de quadrature. 
Nous verrons dans la suite comment obtenir des solutions approchées d'équations différentielles.
Regardons par exemple le problème de Cauchy suivant défini sur l'intervalle $[0,1]$: 
\[
\begin{cases}
    y' = f(t,y) \\
    y(0) = y_0
\end{cases}
\]

Par exemple, ces problèmes regroupent 
\begin{itemize}
    \item l'équation décrivant le lancer d'une balle 
    $k \overrightarrow{v(t)}+m \frac{d \overrightarrow{v(t)}}{d t}=\vec{G}$
    \item l'équation décrivant la charge d'un condensateur 
    $R \frac{d q(t)}{d t}+\frac{1}{C} q(t)=E$
    \item l'équation décrivant la vitesse d'un parachutiste 
    $m \frac{d v(t)}{d t}-g-\frac{1}{2} \rho v(t)^2 S C_x=0$
\end{itemize}

\begin{enumerate}
    \item Parmi les équations ci-dessus, lesquelles sont des équations différentielles linéaires ?
    \ifthenelse{\boolean{showSolutions}}{
        \begin{solution}
            Les deux premières sont linéaires, la troisième non.
        \end{solution}
    }{}
    \item Parmi les équations linéaires, lesquelles sont à coefficients constants ?
    \ifthenelse{\boolean{showSolutions}}{
        \begin{solution}
            Elles sont toutes les deux à coefficients constants.
        \end{solution}
    }{}
    \vspace{1em}
    Reprenons la forme générale du début de l'exercice. 
    L'idée des méthodes numériques que nous utiliserons est de discrétiser l'intervalle $[0,1]$ en $n$ sous-intervalles de longueur $h=1/n$.
    On définit alors 
    \[
    t_i = i h \quad \text{et} \quad y_i = y(t_i)
    \]
    Une solution approchée sera donnée par la valeur de $y_i$ pour $i=0, \cdots, n$.
    \item En intégrant l'équation différentielle sur l'intervalle $[t_i, t_{i+1}]$, déterminer une relation entre $y_{i}$ et $y_{i+1}$.
    \ifthenelse{\boolean{showSolutions}}{
        \begin{solution}
            On intègre l'équation différentielle sur l'intervalle $[t_i, t_{i+1}]$, on obtient : 
            \[
            \int_{t_i}^{t_{i+1}} y'(t) dt = \int_{t_i}^{t_{i+1}} f(t,y(t)) dt
            \]
            On peut alors écrire : 
            \[
        y_{i+1} - y_i = \int_{t_i}^{t_{i+1}} f(t,y(t)) dt
        \]       
        \end{solution}
    }{}
    \item On ne peut pas calculer l'intégrale présente dans l'équation que vous avez trouvé. 
    Proposez $3$ méthodes différentes pour obtenir une valeur approchée de cette intégrale et résoudre l'équation de départ de manière approchée. 
    \ifthenelse{\boolean{showSolutions}}{
        \begin{solution}
            On peut utiliser la méthode des rectangles, la méthode des trapèzes, ou la méthode de Simpson.
        \end{solution}
    }{}
\vspace{1em}
    Chacune de ces méthodes conduira à un algorithme classique de résolution approchée de l'équation différentielle. Nous les détaillerons au prochain TD. 
\end{enumerate}
    

\end{document}
